\section{Disjoint Set (Union-Find)}
\label{sec:graph-disjoint-set-impl}

\problemstatement{Implement a data structure supporting: $\textit{find}(x)$
-- determine which group $x$ currently belongs to -- and
$\textit{union}(x,y)$ -- merge $x$'s and $y$'s groups into one, both
as efficiently as possible.}

\begin{patternbox}
\emph{``Efficiently track and merge groups of connected elements''} is
the textbook use case for Disjoint Set: each element points to a
\textbf{parent}, with a group's identity determined by walking parent
pointers up to a \textbf{root} (a node that is its own parent).
Path compression and union by size, described in
Section~\ref{sec:graph-disjoint-set}, keep both operations nearly
constant time.
\end{patternbox}

\subsection*{Algorithm}

\begin{enumerate}
  \item Initialize: every node is its own parent (its own group), and
        every group has size $1$.
  \item $\textit{find}(x)$: if $x$ is its own parent, return $x$.
        Otherwise, recursively find $x$'s ultimate root, and
        \textbf{set $x$'s parent directly to that root} before
        returning (path compression).
  \item $\textit{union}(x,y)$: find both roots; if equal, they are
        already in the same group, do nothing. Otherwise, attach the
        \textbf{smaller}-size root under the \textbf{larger}-size
        root's parent pointer, and add the smaller group's size into
        the larger's.
\end{enumerate}

\subsection*{C++ implementation}

\begin{lstlisting}
#include <bits/stdc++.h>
using namespace std;

class DisjointSet {
    vector<int> parent, size;
public:
    DisjointSet(int n) : parent(n), size(n, 1) {
        for (int i = 0; i < n; i++) parent[i] = i;
    }

    int find(int x) {
        if (parent[x] == x) return x;
        return parent[x] = find(parent[x]); // path compression
    }

    void unite(int x, int y) {
        int rootX = find(x), rootY = find(y);
        if (rootX == rootY) return;

        if (size[rootX] < size[rootY]) swap(rootX, rootY);
        parent[rootY] = rootX;       // smaller under larger
        size[rootX] += size[rootY];
    }

    bool connected(int x, int y) {
        return find(x) == find(y);
    }
};

int main() {
    DisjointSet ds(6);
    ds.unite(0, 1);
    ds.unite(1, 2);
    ds.unite(3, 4);

    cout << (ds.connected(0, 2) ? "true" : "false") << endl;  // true
    cout << (ds.connected(0, 3) ? "true" : "false") << endl;  // false
    ds.unite(2, 3);
    cout << (ds.connected(0, 4) ? "true" : "false") << endl;  // true
    return 0;
}
\end{lstlisting}

\subsection*{Dry run}

\dryrun{Take $6$ elements, performing $\textit{unite}(0,1)$,
$\textit{unite}(1,2)$, $\textit{unite}(3,4)$, then
$\textit{unite}(2,3)$.}

Initially each node is its own root, size $1$ each.
$\textit{unite}(0,1)$: both size $1$, $\textit{rootX}=0$ stays (tie,
no swap needed in this implementation since $1<1$ is false), parent$[1]=0$,
size$[0]=2$. $\textit{unite}(1,2)$: $\textit{find}(1)=0$ (size $2$),
$\textit{find}(2)=2$ (size $1$); $2$'s group is smaller, parent$[2]=0$,
size$[0]=3$. $\textit{unite}(3,4)$: similarly, parent$[4]=3$,
size$[3]=2$. $\textit{unite}(2,3)$: $\textit{find}(2)=0$ (size $3$,
via path compression directly), $\textit{find}(3)=3$ (size $2$); $3$'s
group is smaller, parent$[3]=0$, size$[0]=5$. Now $0,1,2,3,4$ are all
connected (root $0$); $5$ remains its own separate group.

\begin{complexitybox}
\textbf{Time Complexity.} $O(\alpha(n))$ amortized per
$\textit{find}$/$\textit{union}$ call, where $\alpha$ is the inverse
Ackermann function -- effectively constant for any realistic input
size.
\textbf{Space Complexity.} $O(n)$ for the parent and size arrays.
\end{complexitybox}

\begin{takeawaybox}
Disjoint Set answers ``are these two things in the same group, and can
I merge their groups'' far faster than re-running a DFS/BFS traversal
from scratch after every update -- making it the tool of choice
whenever a problem involves \textbf{repeated} connectivity queries or
incremental merging, rather than a single one-shot traversal.
\end{takeawaybox}
